\documentclass[11pt,a4paper]{article}
\usepackage[italian]{babel}
\usepackage{amsmath,amssymb}
\usepackage[margin=2cm]{geometry}
\usepackage{array}
\usepackage{booktabs}
\usepackage{csquotes}
\renewcommand{\arraystretch}{1.4}
\setlength{\parindent}{0pt}
\setlength{\parskip}{5pt}

\newcommand{\sym}[1]{$#1$}

\begin{document}

\begin{center}
{\Large\bfseries Come si leggono le formule}\\[2pt]
{\itshape Legenda \enquote{alla buona} dei simboli usati nelle risposte d'esame}
\end{center}

Questa guida serve a \textbf{leggere e capire} i simboli che trovi in
\texttt{11\_domande\_esame.pdf}. Non devi impararli a memoria: tienila a fianco
finché non ti vengono naturali.

\section*{1. Lettere greche (sono solo \enquote{nomi})}

Le lettere greche sono solo etichette per quantità: non spaventarti, ognuna ha
un nome e un significato fisso nel corso.

\begin{center}
\begin{tabular}{>{$}c<{$} l l p{7.2cm}}
\toprule
\textrm{Simbolo} & \textrm{Si legge} & \textrm{Nome} & \textrm{Cosa indica nel corso} \\
\midrule
\varphi & \enquote{fi} & phi & funzione di attivazione (ReLU, sigmoide\dots) \\
\psi & \enquote{psi} & psi & funzione di attivazione dell'output \\
\sigma & \enquote{sigma} & sigma & la sigmoide $\sigma(z)=1/(1+e^{-z})$ \\
\theta & \enquote{theta} & theta & i parametri del modello (pesi+bias), o un angolo \\
\eta & \enquote{eta} & eta & learning rate (passo della discesa del gradiente) \\
\delta & \enquote{delta} & delta & il \enquote{segnale di errore} nella backprop \\
\varepsilon & \enquote{epsilon} & epsilon & un numero piccolo a piacere \\
\nabla & \enquote{nabla} & nabla & il gradiente (vettore delle derivate) \\
\Sigma & \enquote{sigma maiusc.} & sigma & sommatoria (somma di tanti termini) \\
\alpha,\ \beta & \enquote{alfa, beta} & alpha, beta & coefficienti generici \\
\gamma & \enquote{gamma} & gamma & esponente della correzione gamma (immagini) \\
\tau & \enquote{tau} & tau & variabile di integrazione (tempo) \\
\bottomrule
\end{tabular}
\end{center}

\section*{2. Apici, pedici e \enquote{cappelli} (la parte che confonde)}

\begin{center}
\begin{tabular}{>{$}c<{$} l p{9.5cm}}
\toprule
\textrm{Scrittura} & \textrm{Si legge} & \textrm{Significato} \\
\midrule
x_i & \enquote{x con i} & l'$i$-esimo elemento del vettore $x$ (pedice = indice) \\
x_0 & \enquote{x zero} & un punto/valore di partenza \\
W^{(1)} & \enquote{W uno} & i pesi del \textbf{layer 1}. L'apice tra parentesi è il \textbf{numero dello strato}, NON un esponente! \\
W^{(2)} & \enquote{W due} & i pesi del layer 2 \\
x^2 & \enquote{x al quadrato} & questo sì è un esponente vero \\
\hat{y} & \enquote{y cappello} & valore \textbf{predetto} dal modello (il \enquote{cappello} = stimato) \\
W^{T} & \enquote{W trasposta} & la matrice ribaltata (righe$\leftrightarrow$colonne) \\
\mathbb{R}^{N} & \enquote{erre enne} & lo spazio dei vettori di $N$ numeri reali \\
\bottomrule
\end{tabular}
\end{center}

Regola d'oro: \textbf{pedice = indice} (\enquote{quale elemento}); \textbf{apice
tra parentesi = numero del layer}; \textbf{apice numerico = potenza}.

\section*{3. Operatori e simboli matematici}

\begin{center}
\begin{tabular}{>{$}c<{$} l p{9.5cm}}
\toprule
\textrm{Simbolo} & \textrm{Si legge} & \textrm{Significato} \\
\midrule
\cdot & \enquote{per} & prodotto (scalare tra vettori: $w\cdot x=\sum_i w_i x_i$) \\
\times & \enquote{per} & moltiplicazione, o dimensioni (es. $224\times 224$) \\
\odot & \enquote{prodotto elemento per elemento} & moltiplica posizione per posizione (Hadamard) \\
* & \enquote{convoluto con} & convoluzione (filtri, CNN) \\
\sum_i & \enquote{sommatoria su i} & somma tutti i termini al variare di $i$ \\
\int & \enquote{integrale} & somma \enquote{continua} (area sotto la curva) \\
\partial & \enquote{derivata parziale} & derivata rispetto a una variabile sola \\
\nabla f & \enquote{gradiente di f} & vettore di tutte le derivate parziali di $f$ \\
\|x\| & \enquote{norma di x} & \enquote{lunghezza} del vettore $x$ \\
|x| & \enquote{valore assoluto} & $x$ senza segno \\
\sqrt{x} & \enquote{radice di x} & radice quadrata \\
\approx & \enquote{circa uguale} & approssimativamente uguale \\
\bottomrule
\end{tabular}
\end{center}

\section*{4. Simboli logici e di insiemi}

\begin{center}
\begin{tabular}{>{$}c<{$} l p{9.5cm}}
\toprule
\textrm{Simbolo} & \textrm{Si legge} & \textrm{Significato} \\
\midrule
\in & \enquote{appartiene a} & $x\in\mathbb{R}$: \enquote{x è un numero reale} \\
\forall & \enquote{per ogni} & vale per tutti \\
\exists & \enquote{esiste} & esiste almeno uno \\
\Rightarrow & \enquote{implica} & se\dots allora \\
\Leftrightarrow & \enquote{se e solo se} & vale in entrambi i versi \\
\to & \enquote{tende a / va in} & $f:\mathbb{R}^N\to\mathbb{R}$: prende $N$ numeri, restituisce 1 numero \\
\le,\ \ge & \enquote{minore/maggiore uguale} & disuguaglianze \\
\bottomrule
\end{tabular}
\end{center}

\newpage
\section*{5. Le formule chiave \enquote{lette a parole}}

\subsection*{Rete neurale (MLP con un layer nascosto)}
\[ h = \varphi(W^{(1)} x + b^{(1)}), \qquad \hat{y} = \psi(W^{(2)} h + b^{(2)}) \]
\textbf{Si legge}: \enquote{acca uguale fi di W-uno per x più b-uno; y-cappello
uguale psi di W-due per acca più b-due}.\\
\textbf{Significato}: prendo l'input $x$, lo combino con i pesi $W^{(1)}$, aggiungo
il bias $b^{(1)}$ e applico l'attivazione $\varphi$ $\rightarrow$ ottengo lo strato
nascosto $h$. Poi rifaccio lo stesso con $h$ per ottenere la predizione $\hat{y}$.
È \textbf{due volte lo stesso schema} \enquote{pesi$\times$input + bias, poi attivazione}.

\subsection*{Discesa del gradiente}
\[ x_{k+1} = x_k - \eta \cdot \nabla f(x_k) \]
\textbf{Si legge}: \enquote{x al passo k+1 uguale x al passo k meno eta per il
gradiente di f}.\\
\textbf{Significato}: dal punto attuale mi sposto un pochino ($\eta$) nella
direzione opposta al gradiente (cioè in \enquote{discesa}), per avvicinarmi al minimo.

\subsection*{Funzione di costo MSE}
\[ L(\theta) = \frac{1}{n}\sum_i (f_\theta(x_i) - t_i)^2 \]
\textbf{Si legge}: \enquote{L di theta uguale uno su n, sommatoria su i, di
(predetto meno vero) al quadrato}.\\
\textbf{Significato}: media degli errori al quadrato tra predizione $f_\theta(x_i)$
e valore vero $t_i$. Più è piccola, meglio il modello.

\subsection*{Filtro di Sobel (bordi)}
\[ M = \sqrt{G_x^2 + G_y^2}, \qquad \theta = \arctan2(G_y, G_x) \]
\textbf{Si legge}: \enquote{M uguale radice di G-x quadro più G-y quadro}.\\
\textbf{Significato}: $G_x,G_y$ sono le variazioni di luminosità in orizzontale e
verticale; $M$ è quanto è \enquote{forte} il bordo, $\theta$ la sua direzione.

\subsection*{Softmax + cross-entropy (classificazione)}
\[ \text{softmax}(z)_i = \frac{e^{z_i}}{\sum_j e^{z_j}}, \qquad L = -\log p_c \]
\textbf{Si legge}: \enquote{softmax di z, componente i, uguale e alla z-i diviso la
somma di e alla z-j}.\\
\textbf{Significato}: la softmax trasforma dei punteggi in \textbf{probabilità}
(che sommano a 1); la cross-entropy $-\log p_c$ punisce il modello quanto più
bassa è la probabilità data alla classe giusta $c$.

\vspace{1em}
\hrule
\vspace{6pt}
\textbf{Trucco per leggere qualsiasi formula}: leggi da sinistra a destra dicendo
ad alta voce cosa fa ogni pezzo (\enquote{prendo questo, lo moltiplico per\dots,
aggiungo\dots, applico\dots}). Se sai \emph{a parole} cosa fa, la formula la
\textbf{ricostruisci} da solo: all'esame non serve ricordarla a memoria, serve
capirla.

\end{document}
